\subsubsection{The Strong Maximum Principle}\label{subsec:parabolic-strong-maximum}

The goal of this section is the following result.

\begin{Theorem}
\label{thm:parabolic-strong-maximum}
Assume that $Lu\geq0$ ($Lu\leq0$). Let $M=\sup_Du$ ($M=\inf_Du$). Assume that $u=M$ at a point $(\bm{x}_0,t_0)\in D$ and that one of the following holds:
\begin{enumerate}
\item $c=0$ and $M$ is arbitrary,
\item $c\leq0$ and $M\geq0$ ($M\leq0$),
\item $M=0$ and $c$ is arbitrary.
\end{enumerate}
Then $u=M$ on $\overline\Omega\times[0,t_0]$.
\end{Theorem}

The proof follows from three lemmas. From now on, we always assume $Lu\geq0$ without explicitly stating so; the case $Lu\leq0$ obviously follows by substituting $-u$ for $u$. The following lemma is immediate from Lemma~\ref{lem:elliptic-boundary-point} and Remark~\ref{rem:boundary-point-semidefinite}.

\begin{Lemma}
\label{lem:parabolic-ball-extreme-time}
Let $B\subset\mathbb{R}^{n+1}$ be a ball with $\overline B\subset D$ and suppose $u<M$ in $B$ and $u(\bm{x}_0,t_0)=M$, where $(\bm{x}_0,t_0)\in\partial B$. Then $t_0$ is either the smallest or the largest value which $t$ assumes in $\overline B$.
\end{Lemma}

Next we show that if $u<M$ at any point $(\bm{x}_0,t_0)$ in $Q$, then $u<M$ on all of $\Omega\times\{t_0\}$.

\begin{Lemma}
\label{lem:parabolic-slice-strictness}
Assume that $u(\bm{x}_0,t_0)<M$, where $\bm{x}_0\in\Omega$ and $t_0\in(0,T)$. Then $u(\bm{x},t_0)<M$ for every $\bm{x}\in\Omega$.
\end{Lemma}

\begin{Proof}
Assume that $\Omega\times\{t_0\}$ contains points where $u=M$. Then we can find points $\bm{x}_1$ and $\bm{x}_2$ in $\Omega$ such that $u(\bm{x}_1,t_0)=M$, $u(\bm{x}_2,t_0)<M$, and the line segment $L$ connecting $(\bm{x}_1,t_0)$ and $(\bm{x}_2,t_0)$ lies entirely in $\Omega\times\{t_0\}$. By moving the point $\bm{x}_1$, we can always achieve that $u<M$ along $L$. Now let
\begin{equation}\label{eq:parabolic-slice-distance-radius}
\delta=\min(|\bm{x}_1-\bm{x}_2|,\operatorname{dist}(L,\partial D)).
\end{equation}
For $\bm{x}\in L$ with $0<|\bm{x}-\bm{x}_1|<\delta$, let
\begin{equation}\label{eq:parabolic-slice-maximum-set-distance}
d(\bm{x})=\operatorname{dist}((\bm{x},t_0),Q\cap\{u=M\}).
\end{equation}
Obviously, $d(\bm{x})\leq|\bm{x}-\bm{x}_1|$.

By Lemma~\ref{lem:parabolic-ball-extreme-time}, the point in $Q\cap\{u=M\}$ which is nearest to $(\bm{x},t_0)$ is of the form $(\bm{x},t)$, so that either $u(\bm{x},t_0+d(\bm{x}))=M$ or $u(\bm{x},t_0-d(\bm{x}))=M$. Let $\bm{n}$ be a unit vector in the direction of $L$. For sufficiently small $|\epsilon|>0$, $d(\bm{x}+\epsilon\bm{n})$ is defined, and by the Pythagorean Theorem we have
\begin{equation}\label{eq:parabolic-slice-distance-forward-bound}
d(\bm{x}+\epsilon\bm{n})\leq\sqrt{\epsilon^2+d(\bm{x})^2}.
\end{equation}
By the same argument,
\begin{equation}\label{eq:parabolic-slice-distance-reverse-bound}
d(\bm{x})\leq\sqrt{\epsilon^2+d(\bm{x}+\epsilon\bm{n})^2},
\end{equation}
and hence
\begin{equation}\label{eq:parabolic-slice-distance-lower-bound}
d(\bm{x}+\epsilon\bm{n})\geq\sqrt{-\epsilon^2+d(\bm{x})^2}.
\end{equation}
From \eqref{eq:parabolic-slice-distance-reverse-bound} and \eqref{eq:parabolic-slice-distance-lower-bound}, it follows that the derivative of $d(\bm{x}+\epsilon\bm{n})$ at $\epsilon=0$ exists and is zero. Hence $d(\bm{x})$ is constant along $L\cap\{0<|\bm{x}-\bm{x}_1|<\delta\}$. But this is a contradiction since $d(\bm{x})\ne0$, but $d(\bm{x})\to0$ as $\bm{x}\to\bm{x}_1$.
\end{Proof}

The final lemma which we require is the following.

\begin{Lemma}
\label{lem:parabolic-time-strictness}
Let $0\leq t_0<t_1\leq T$ and assume that $u<M$ in $\Omega\times(t_0,t_1)$. Then $u<M$ on $\Omega\times\{t_1\}$.
\end{Lemma}

\begin{Proof}
Assume on the contrary that there is $\bm{x}_1\in\Omega$ with $u(\bm{x}_1,t_1)=M$. Define
\begin{equation}\label{eq:parabolic-time-propagation-barrier}
v(\bm{x},t)=\exp\bigl(-|\bm{x}-\bm{x}_1|^2-\alpha(t-t_1)\bigr)-1,
\end{equation}
where $\alpha>0$ is chosen large. We compute\[
\begin{aligned}
Lv(\bm{x},t)&=\exp\bigl(-|\bm{x}-\bm{x}_1|^2-\alpha(t-t_1)\bigr)\\
&\quad\times\bigl(4a_{ij}(x_i-x_{1i})(x_j-x_{1j})-2(a_{ii}+b_i(x_i-x_{1i}))+\alpha\bigr)+cv.
\end{aligned}
\]If $N$ is a small neighborhood of $(\bm{x}_1,t_1)$, we obtain $Lv>0$ in $N\cap\{t\leq t_1\}$ if we choose $\alpha$ large enough. Let now $A$ be the domain bounded by the paraboloid $|\bm{x}-\bm{x}_1|^2+\alpha(t-t_1)=0$ and a sufficiently small sphere centered at $(\bm{x}_1,t_1)$. The boundary of $A$ has two parts, one on the paraboloid, where $v=0$, and one on the sphere and inside the paraboloid, where $u<M$.

If we choose $\epsilon>0$ small, we therefore have $u+\epsilon v-M\leq0$ on $\partial A$ and $L(u+\epsilon v-M)\geq-cM\geq0$ in $A$.

If $c\leq0$, the weak maximum principle implies that $u+\epsilon v-M\leq0$ throughout $A$. Taking the $t$ derivative at the point $(\bm{x}_1,t_1)$, we conclude
\begin{equation}\label{eq:parabolic-terminal-time-derivative-bound}
u_t+\epsilon v_t=u_t-\epsilon\alpha\geq0,
\end{equation}
i.e., $u_t>0$. But since $(\bm{x}_1,t_1)$ is a maximum with respect to $\bm{x}$, we find that at this point
\begin{equation}\label{eq:parabolic-terminal-operator-bound}
Lu\leq-u_t+cu\leq-u_t,
\end{equation}
and since $Lu\geq0$ was assumed, we conclude $u_t\leq0$, a contradiction.

If $M=0$, we can, as in the proof of Lemma~\ref{lem:elliptic-boundary-point}, apply the same argument with $(L-c^+)$ in place of $L$.
\end{Proof}

The proof of the theorem is now clear. For every $t>0$, we have either $u(\bm{x},t)=M$ or $u(\bm{x},t)<M$ for all $\bm{x}\in\Omega$ (Lemma~\ref{lem:parabolic-slice-strictness}). The set $Z$ of all $t$\textquotesingle s for which $u<M$ is open and hence a countable union of intervals. By Lemma~\ref{lem:parabolic-time-strictness} none of these intervals can have an upper endpoint. The only possibility is then that $Z=(t_0,T]$ for some $t_0$, which is the claim of the theorem.

Finally we remark that if $\partial\Omega$ is of class $C^2$ and $u$ is $C^1$ up to the boundary, then Lemma~\ref{lem:elliptic-boundary-point} applies, i.e., if $\bm{x}_0\in\partial\Omega$, $t\in(0,T)$, $u(\bm{x}_0,t)=M$, but $u(\bm{x},t)<M$ for $\bm{x}\in\Omega$, then $\partial u/\partial n(\bm{x}_0,t)>0$.

\paragraph*{Problems}

\begin{enumerate}[label=\thesection.\arabic*.,ref=\thesection.\arabic*,start=22,series=book-problems]

\item\label{ex:parabolic-neumann-comparison} Assume that $\Omega$ is bounded, $\partial\Omega$ is of class $C^2$ and that $u,v\in C^2(D)\cap C^1(\overline D)$. Assume, moreover, that $Lu\leq Lv$ for $(\bm{x},t)\in D$, that $u(\bm{x},0)\geq v(\bm{x},0)$ for $\bm{x}\in\Omega$ and that $\partial u/\partial n\geq\partial v/\partial n$ for $(\bm{x},t)\in\partial\Omega\times(0,T)$. Show that $u\geq v$ in $D$.

\item\label{ex:parabolic-stable-equilibrium-nonnegative} Assume that $\Omega$ is bounded, $f\leq0$ and that $f$ and the coefficients of $L$ are independent of $t$. We consider the equation $Lu=f$ with boundary conditions $u=0$ on $\partial\Omega\times(0,\infty)$. An equilibrium solution is a solution which does not depend on time. An equilibrium solution $v$ is called stable if, for every continuous function $u_0(\bm{x})$ satisfying $u_0=0$ on $\partial\Omega$, there is a solution $u$ satisfying $Lu=f$ in $\Omega\times(0,\infty)$, $u=0$ on $\partial\Omega\times(0,\infty)$ and $u(\bm{x},0)=u_0(\bm{x})$ for $\bm{x}\in\Omega$ and, moreover, $u(\bm{x},t)\to v(\bm{x})$ as $t\to\infty$. Show that a stable equilibrium solution must be non-negative.

\item\label{ex:whole-space-bounded-heat-maximum} Let $D=\mathbb{R}^n\times(0,T]$. Let $u\in C^2(D)\cap C(\overline D)$ be a bounded solution of the heat equation $u_t=\Delta u$. Show that $\sup_Du\leq\sup_{\mathbb{R}^n}u(\cdot,0)$. Hint: Consider the function $v=u-\epsilon(2nt+|\bm{x}|^2)$.

\item\label{ex:semilinear-parabolic-uniqueness} Let $\Omega$ be bounded, and let $f:\mathbb{R}\to\mathbb{R}$ be of class $C^1$. Let $u_0\in C(\overline\Omega)$ be such that $u_0=0$ on $\partial\Omega$. Prove that the equation $u_t=\Delta u+f(u)$, subject to the boundary condition $u=0$ on $\partial\Omega$ and the initial condition $u(\bm{x},0)=u_0(\bm{x})$, has at most one solution on $\Omega\times(0,T)$, for any $T>0$.

\item\label{ex:reaction-diffusion-positivity} Let $\Omega$ be bounded and let $f,g:\mathbb{R}^2\to\mathbb{R}$ be of class $C^1$. Let $u$, $v$ satisfy the equations $u_t=\Delta u+f(u,v)$, $v_t=\Delta v+g(u,v)$ on $\Omega\times(0,T)$, with boundary conditions $u=v=0$ on $\partial\Omega$. Assume that $f(0,v)\geq0$ for every $v$ and that $g(u,0)\geq0$ for every $u$. Show that if $u(\bm{x},0)\geq0$, $v(\bm{x},0)\geq0$ for every $\bm{x}\in\Omega$, then $u\geq0$, $v\geq0$. Systems of the kind described here arise in reaction-diffusion problems and in mathematical biology, where $u$ and $v$ denote concentrations or populations of species.

\end{enumerate}
